% Readout SAE diagnostics and reproducibility appendix.

\section{Readout SAE Diagnostics and Reproducibility}
\label{app:readout-factorizer-diagnostics-reproducibility}

This appendix collects secondary diagnostics, alternative-recipe comparisons,
denominator controls, and reproducibility details for the selected readout SAEs.

\subsection{Additional Metric Analyses}
\label{app:readout-factorizer-diagnostics}

These analyses ask whether the remaining error comes from coefficient fitting,
support selection, low-rank structure, or reconstruction-vs-readout metric
mismatch.

\paragraph{Support selection.}
On the Qwen3.5-9B $D=32768$, $k=128$ recipe, the trained encoder reaches rowEV
$0.621$. Nonnegative least-squares \citep{lawson1974solving} on the encoder's
own support reaches $0.625$, nonnegative OMP \citep{pati1993orthogonal}
$0.647$, and signed-OMP $0.655$. The coefficient gap is therefore small
(about $0.004$), while same-coefficient-class support-selection headroom is
about $0.026$. Dense rank-128 reconstruction reaches only $0.101$, far below the sparse
dictionary, ruling out low-rank PCA structure
\citep{jolliffe2002principal}.
Figure~\ref{fig:app-k-omp-diagnostic} shows both the strict-budget recipe and
the high-fidelity 32$\times$, $k=256$ run. The right-hand table labels the
latter OMP numbers as sparse-reconstruction ceilings; deployed-encoder metrics
are the trained-encoder rows.

\begin{figure*}[!tbp]
    \centering
    \begin{minipage}[t]{0.60\textwidth}
        \vspace{0pt}
        \centering
        \ReadoutSAEGraphic[width=\linewidth,height=0.27\textheight]{appendix_k_omp_diagnostic.pdf}
    \end{minipage}\hfill
    \begin{minipage}[t]{0.35\textwidth}
        \vspace{0pt}
        \centering
        \scriptsize
        \setlength{\tabcolsep}{4pt}
        \renewcommand{\arraystretch}{1.08}
        \textbf{32$\times$, $k=256$ ceiling comparison}\par\vspace{2pt}
        \begin{tabular}{@{}lcc@{}}
        \toprule
        Quantity & rowEV & 95\% CI \\
        \midrule
        Trained encoder & $0.858$ & $[0.856, 0.859]$ \\
        Signed LS support & $0.864$ & $[0.863, 0.866]$ \\
        Signed-OMP proxy & $0.902$ & $[0.900, 0.903]$ \\
        \bottomrule
        \end{tabular}
    \end{minipage}
    \caption{
        \textbf{Support-selection analyses for the Qwen3.5-9B readout SAE.}
        Left: the strict-budget $k=128$ recipe and the high-fidelity
        32$\times$, $k=256$ capacity run. Right: 32$\times$, $k=256$ EV-ceiling
        numbers with 95\% bootstrap CIs. OMP (orthogonal matching pursuit) is
        reported as an achieved sparse-reconstruction comparison.
    }
    \label{fig:app-k-omp-diagnostic}
\end{figure*}

For the 32$\times$, $k=256$ capacity run, coefficient refitting contributes
$+0.006$ rowEV, while support selection contributes $+0.038$. The encoder
reaches $95.1\%$ of the signed-OMP proxy, similar to the $D=32768$, $k=128$
ratio. Dense rank-256 reconstruction is $0.155$, so the $k=256$ sparse code is
about $5.9\times$ better by unexplained-variance reduction than the same-rank
dense reference, since \((1-0.155)/(1-0.858)\approx5.9\).

\paragraph{Small-model support-selection analysis.}
We repeat the analysis on the 16$\times$, $k=128$ small-model transfer 5k
checkpoints to isolate the source of the remaining reconstruction gap with the
learned dictionary fixed; the converged transfer metrics are reported in
Appendix~\ref{app:readout-factorizer-small-model-transfer}. Across all four models, refitting
coefficients on the encoder's own support gives a much smaller gain than
changing the support, indicating that support selection is the dominant
remaining recipe-side bottleneck. The signed-OMP gaps are $+0.158$ for
Qwen3.5-0.8B, $+0.099$ for Qwen3.5-2B, $+0.182$ for Gemma-4-E2B, and
$+0.125$ for Gemma-4-E4B.

\ReadoutSAEFloatBarrier

\paragraph{Metric mismatch.}
The capacity sweep motivates selecting readout SAEs with both row-level and
readout-level metrics. The $D=65536$, $k=256$ model truncated to $k=128$ has
higher rowEV than the fixed-$k=128$ recipe ($0.695>0.621$), while top1/KL favor
the fixed-$k=128$ recipe ($0.721/0.469$ vs $0.846/0.296$). The $D=131072$
truncation shows the same metric separation: rowEV/top1/KL are
$0.816/0.738/0.425$. These truncated comparisons motivate the selection rule:
use rowEV, top1, KL, reconstruction-error metrics, and usage statistics jointly.

\ReadoutSAEFloatBarrier
\subsection{Alternatives and Negative Controls}
\label{app:readout-factorizer-controls}

\paragraph{Alternative recipes and denominator controls.}
Matched comparisons rule out the non-selected recipes.
At $D=32768$/5k, plain TopK exceeded the matryoshka variants in rowEV
\citep{bussmann2025matryoshka} ($0.567$ vs.\ $0.559/0.554$) at lower cost, and
the $D=16384$ fixed-$k=128$ cell trailed $D=32768$ by about $0.06$ rowEV.
Calibrated BatchTopK \citep{bussmann2024batchtopk} reached rowEV $0.531$ on the
matched frontier cell. Qwen JumpReLU \citep{rajamanoharan2024jumprelu} missed
the intended target $L_0=256$ (achieved about $38$--$68$), and step- or
sparsity-matched TopK controls had higher rowEV ($0.42$--$0.44$ vs.\
$0.26$--$0.31$), with mixed logit metrics. Finally, tokenizer/non-token-row
audits support the denominator choice: structural non-token rows were $5.4\%$ of
rows, and excluding non-token readout rows changed rowEV by $+0.0016$.

\subsection{Reproducibility and Selection Details}
\label{app:readout-factorizer-limitations}

The main experiments report exact-versus-reconstructed scores, sign agreement,
relative error, and null/reference baselines for interpreted scores.
Appendices~\ref{app:model-suite-selection}--\ref{app:robustness-controls}
specify models, dictionary settings, seeds, row sampling, selection order,
tokenization filters, score construction, and denominator controls. Dictionary
selection uses row reconstruction, replacement fidelity, and sparse code usage;
we reserve held-out logit differences for later checks. Remaining
risks: computational cost, checkpoint/tokenizer versioning, and
tokenizer-specific filtering for single-token scores. Local decompositions include
explicit reconstruction errors, and we restrict qualitative interpretations
to sufficiently reconstructed scores.

We selected the Qwen/Gemma readout SAEs before downstream qualitative
interpretation, using row reconstruction, replacement-logit metrics, and
sparse code usage statistics. We reserved the logit-difference contrast set
for the later reconstruction checks, so selected-score reconstruction is
evaluated after dictionary selection. Most reported dictionary settings use seed \(0\), with the
Qwen3.5-0.8B 32$\times$, $k=256$ three-seed window providing the
seed-variation snapshot above.
For Qwen3.5-9B, tokenizer-free denominator controls replace tokenizer-aligned
text-token-only EV and frequency-stratified fidelity: the tokenizer row
count did not cleanly align with the $\WU$ rows, and no vocab-aligned unigram
table was available. The tokenizer-free non-token-row audit therefore serves as
the denominator-control result for that model. Gemma rowEV is measured on the
linear pre-softcap unembedding rows and is cap-independent; Gemma top1/KL are
evaluated under the model's $30\tanh(\cdot/30)$
\texttt{final\_logit\_softcapping}. Since the cap is strictly monotone, top1 is
cap-independent and the cap affects only KL.
The selected operating points and reported claims come from the Qwen/Gemma
sweeps and the targeted Ministral/R1 comparison rows above; earlier Pythia-160M
architecture and continuity sweeps \citep{biderman2023pythia} are
design-history records.
Exact replication depends on using the same model checkpoint revisions,
tokenizers, final-normalization and softcap code paths, score filters, and
single-token eligibility rules. For this reason, every interpreted local score
decomposition reports the exact score, reconstructed score, reconstruction error, and,
for signed contrasts, sign agreement.

\paragraph{Training and evaluation recipe.}
\begingroup
\footnotesize
\sloppy
All reported readout SAEs are TopK sparse autoencoders trained on the
model's final readout rows. The reported Qwen/Gemma recipe uses
row-seeded initialization, \texttt{hybrid\_50freq\_50uniform} row sampling, a
warmup-cosine learning-rate schedule from \(10^{-3}\) to \(10^{-4}\), and seed
\(0\) unless stated otherwise. The Qwen3.5-9B recipe-selection run uses
\(D=32768\), native \(k=128\), a delayed prism-loss ramp to \(10^{-3}\) after
\(70\%\) of training, and 20k optimization steps; the capacity and transfer
runs keep the same recipe family while varying \(D\) and native \(k\) as
reported in Tables~\ref{tab:app-k-model-finalists},
\ref{tab:app-k-qwen-capacity}, and
\ref{tab:app-k-small-model-converged-ledger}. Hidden-state evaluation uses the
final readout after the relevant final normalization and, for Gemma, the
model's monotone final-logit softcap for top1/KL evaluation. The reported
metrics are rowEV, top-1 agreement, KL in bits, sparse code usage, and the
selected-score reconstruction error and sign agreement defined above.
\endgroup

\ReadoutSAEFloatBarrier
